\chapter{Methodology}
\label{chap:method}

\section{Research Design}

The quantitative and experimental research design is followed in the present study because it will
develop and test machine learning to forecast the post-harvest loss risk. A quantitative research
method is selected because the given study aims at examining structured numerical data and
determining relationships between the independent variables and the target outcome variable of the
research. The experimental aspect involves the training and testing of various machine learning
models to determine their predictive performance.

This design is appropriate because the research questions are objective in nature, measuring the
relationship between factors that include storage conditions, transportation, and environmental
variables and their effect on loss after harvest. The quantitative approach will make it possible
to statistically assess results and compare individual models on objective grounds. Moreover, the
experimental methodology facilitates reproducibility because models can be re-trained on exactly
the same dataset and parameters. This design is hence suitable in the creation of a sound,
data-driven decision support instrument for smallholder agribusinesses.


\section{Research Onion Explanation}

The research is grounded in the research onion model, which is conceptualized as a series of
methodological choices layered from the outside in.

At the philosophical level, the research is based on \textbf{positivism}, as the researcher assumes
that reality can be measured objectively using observable data. This is suitable because the
variables included are measurable, such as temperature, conditions of storage, and time spent on
transport.

The study uses a \textbf{deductive approach}, starting with the existing theory on post-harvest
losses and machine learning, which is then tested using empirical data analysis. Hypotheses on the
effect of certain factors and model performance are tested using statistical procedures.

The study design is \textbf{experimental} in the sense that the results of various machine learning
algorithms are compared upon the same dataset to assess predictive accuracy.

The \textbf{time horizon} is cross-sectional, where data are analyzed at a particular point in
time rather than over long periods. This is appropriate given the emphasis on forecasting loss risk
depending on present conditions, as opposed to studying long-term trends.


\section{Data Collection}

This research uses secondary data as the main information source. Specifically, the data is
sourced from the Nigeria General Household Survey (GHS-Panel Wave~5, 2023--2024)~\cite{ghspanel2024nigeria},
which is part of the World Bank's Living Standards Measurement Study -- Integrated Surveys on
Agriculture (LSMS-ISA). This survey provides organizational and country-based data concerning
smallholder agriculture in Nigeria.

The data is organized and table-like, and hence machine learning applications can readily be
applied to it. It includes both numerical and categorical variables involved in the post-harvest
processes. The key features obtained include:

\begin{itemize}
    \item Storage conditions (e.g.\ traditional versus improved storage methods)
    \item Time and distance of transportation
    \item Environmental factors (e.g.\ rainfall and temperature)
    \item Management of practices in harvest and post-harvest phases
\end{itemize}

The post-harvest loss risk target variable is calculated in accordance with available data based on
the comparison of harvested and sold quantities, allowing each record to be classified as either
high-loss or low-loss.


\section{Data Preprocessing}

Preprocessing of data is a very important measure to guarantee the size and quality of the data.
The process begins with \textbf{data cleaning}, at which point inconsistencies, duplicates, and
irrelevant variables are eliminated.

\textbf{Missing values} are handled using appropriate techniques such as imputation or elimination,
based on the quantity and significance of the missing data. \textbf{Categorical variables} (e.g.\
type of storage and handling practises) are converted into numerical representation using encoding
schemes like one-hot encoding or label encoding, to ensure compatibility with machine learning
algorithms.

\textbf{Numerical variables}, wherever applicable, are standardized or normalized by adopting equal
scaling of features. This is especially important for models that are sensitive to the magnitude of
features, such as Logistic Regression. Preprocessing in general guarantees the appropriate
organization, consistency, and aptness of the dataset for proper model training and evaluation.


\section{Machine Learning Models}

This study uses three machine learning models: Logistic Regression, Decision Trees, and Random
Forests. Logistic Regression is an easy-to-understand model adopted as the base model. It works
well when dealing with binary classification issues and gives explicit information on how the
independent variables relate to the target outcome.

Decision Trees are included because of their rule-based structure, which makes the interpretation of
decision paths easy. They are able to find non-linear correlations and can deal with both nominal
and numerical data.

Random Forest is employed as an ensemble learning technique to enhance predictive performance. It
minimizes overfitting by utilizing many decision trees and augments generalization. It is applicable
especially when identifying intricate interactions among variables.

These models have been selected because of their simplicity, interpretability, and ability to work
with small to medium-sized datasets. They are less resource-intensive and easier to apply in a
practical smallholder setting than deep learning models.


\section{Model Evaluation Metrics}

The models are tested based on four primary metrics: accuracy, precision, recall, and F1-score,
supplemented by AUC-ROC.

\textbf{Accuracy} considers the total correctness of predictions. It can however be misleading when
an imbalanced dataset is used. \textbf{Precision} measures the fraction of correctly predicted
positive cases in reference to the total number of predicted positives. The importance of high
precision lies in minimizing false alarms. \textbf{Recall} determines the model's capability to
correctly identify actual positive cases. This is especially pertinent in this research as failure
to identify high-risk situations of loss would result in poor decision-making. The \textbf{F1-score}
balances precision and recall and is a good measure to use when data is imbalanced. The
\textbf{AUC-ROC} (Area Under the Receiver Operating Characteristic Curve) measures the overall
discriminative ability of the model across all classification thresholds. All of these metrics
together give a complete analysis of model performance.


\section{Feature Importance Analysis}

An analysis of the influence of features on post-harvest loss risk is performed to identify the most
influential factors. This is primarily accomplished through the Random Forest model, which provides
a ranking of features according to their contribution to prediction accuracy. The analysis of
feature importance allows the study to identify crucial drivers such as storage conditions, time of
transportation, and environmental factors. This not only enhances the interpretability of the model
but is also useful in making convenient decisions that show areas where interventions are likely to
be most effective.

\section{Code Repository}

The complete source code, data preprocessing scripts, and model training pipelines developed
for this project are publicly available in the following GitHub repository:

\begin{center}
    \url{https://github.com/achauque14/ML-Post-harvest-loss-risk-project}
\end{center}

The repository contains all Python scripts used for data cleaning, model training, evaluation,
and visualization, along with instructions for reproducing the results presented in this
dissertation.


\section{Ethical Considerations}

Ethical issues are put into consideration during the research process. The research involves the
use of secondary data found in publicly accessible materials, and as a result, there is adherence
to policies governing the use of such data. No personal identity information is released. Data
privacy and confidentiality are ensured, and the analysis is conducted in a responsible manner in
order to prevent improper use or misleading results. Also, the research supports the responsible
use of artificial intelligence, concentrating on the transparency and impartiality of model
development.


\section{Limitations of Methodology}

This research is also methodologically limited in a number of aspects. Although the dataset used
was adequate for the analysis, its size might limit the generalisability of findings. There is
potential data bias (reporting error or regional bias) that may influence model performance. In
addition, the use of cross-sectional data limits the temporal analysis of post-harvest losses.
Lastly, the emphasis on simple models that are good in interpretability can lead to worse
predictive power than more complex methods.
